% ============================================================
% sections/s04_experiments.tex —— 4 仿真实验与结果分析
% 写作要点：
%   4.1 实验设置：想定场景、参数取值表、硬件/软件环境、评价指标定义
%   4.2 对比方法与消融设计：对比算法（2~4 个 baseline）+ 本文方法消融变体
%   4.3 结果与分析：主实验表/图 + 定量分析（提升百分比、收敛曲线、
%       方案甘特图/分配结果示意图）
%   4.4 讨论（可选）：敏感性分析（参数影响）、可扩展性（规模增长）
% 注意：实验口径与第 2 节模型约束一一对应；统计显著性/多次独立运行均值±方差。
% ============================================================

\section{仿真实验与结果分析}

\subsection{实验设置}

\placeholder{场景与参数：任务规模、无人机规模、约束参数取值表（表 4.1，三线表）；
运行环境（CPU/内存/语言/求解器）；每组实验独立运行 30 次取均值等口径说明。}

\placeholder{评价指标定义：目标函数值、可行率、求解时间、均衡度等（公式 (4.x)）。}

\subsection{对比方法}

\placeholder{对比算法列表与来源文献（如：遗传算法、差分进化、粒子群、
商业求解器精确解等）；本文方法的消融变体设置（去掉改进点 ①/② 后的版本）。}

\subsection{结果与分析}

\placeholder{主结果表（表 4.2）：各方法在不同规模下的指标对比，
标注最优/次优；定量结论（本文方法较对比方法提升 xx\%）。}

\placeholder{图 4.x：收敛曲线对比；图 4.x：典型场景任务分配结果示意
（无人机-任务连线图或甘特图）。逐图给出分析文字。}

\subsection{敏感性分析}

\placeholder{（可选）关键参数（种群规模、惩罚系数、权重系数等）
对结果的影响曲线 + 结论。}
